% a06_BALtorus_A1: part of Step 3 of the proof of thm:a06:T1 (input from sec/a06_BALtorus_c.tex, after the part files d1-ct8, before Step 4).
% Group (A), part 1 of 4: setting, cube coercivity, coercivity with mixed pins, the Combes-Thomas bound.
% Continues in sec/a06_BALtorus_A2.tex (block inversion, resolvent differences, (1.18)), A3 (depth bounds, (5.5), (1.20), Ba1 Prop. 2.2, the summary theorem), A4 (the torus).

\subsubsection*{Step 3(A): the propagator estimates of group (A)}

This part proves the estimates of group (A) that Steps~3(B)--3(D) consume: the coercivity \cite[(1.8)]{Ba4} on a general region and for the mixed-pin operators \cite[(3.5)]{Ba4}, the Combes--Thomas form of \cite[Cor.~2.3]{Ba4} for every pin structure, the two-sided bounds \cite[(1.15)]{Ba4}, the decay \cite[(1.16)]{Ba4}, the conditioning-region difference \cite[(1.18)]{Ba4}, the domain differences \cite[(5.5), (1.20)]{Ba4}, \cite[Prop.~2.2, (2.27)--(2.29)]{Ba1}, the two-sided form of the depth factor of \cite[Thm., (1.12)]{Ba4}, and their validity on the torus with torus-free constants. The notation is that of Steps~3(B)--3(D): $\eta=L^{-k}$, $e_k=et_k^{1/2}$, $a_k\in[a_*,a]$ with $a_*=a(1-L^{-2})$, $Q_k(A)$, $P_k(A)=Q_k^*Q_k$, $G_k(\Omega,A)$, $\Delta^{(k)}(\Omega,A)$, $C^{(k)}_\Lambda=C_\Lambda$, and $M$ the large-block side. Distances are the sup-distances of \cite[(1.3)]{Ba1} (on a torus, the quotient distances); $|x-x'|$ is the sup-distance of unit-lattice points, and kernels of $N$-component operators are $N\times N$ matrices, $|K(x,x')|$ their operator norm.

\paragraph{Setting.} Functions on the fine lattice $\eta\Z^d$ carry the $\eta^d$-weighted $\ell^2$ product $\langle\phi,\psi\rangle=\sum_x\eta^d\phi(x)\cdot\psi(x)$, functions on the unit lattice the unweighted one; bond functions are weighted at the initial point $b_-$ of the bond. $U(A_b)\defeq\exp(qe_k\eta A_b)$ with $q$ a fixed real antisymmetric $N\times N$ matrix \cite[(1.2)]{Ba4}; all $U(A_b)$ commute and are orthogonal, and $U(A(\Gamma))\defeq\prod_{b\in\Gamma}U(A_b)$ for a contour $\Gamma$. For a union $\Omega$ of unit blocks, $(D_A\phi)(b)\defeq\eta^{-1}(U(A_b)\phi(b_+)-\phi(b_-))$ on the bonds $b\subset\Omega$, so that $\langle\phi,-\Delta^{\eta,N}_{A,\Omega}\phi\rangle=\norm{D_A\phi}^2$ is the Neumann covariant Dirichlet form \cite[(1.3)]{Ba4}. $Q_k(A)$ averages over the unit blocks $B^k(y)$, $y\in\Omega^{(k)}$, along the contours $\Gamma^{(k)}_{y,x}\subset B^k(y)$ of \cite[(2.1)--(2.2)]{Ba1}; $q_k(y)\defeq Q_k(A)^*\delta_y$ is supported in $B^k(y)$, $|q_k(y)(x)|=1$ there, $\norm{q_k(y)}=1$; $P_k(A)=Q_k^*Q_k$ is the orthogonal projection onto the transported constants $x\mapsto U(A(\Gamma^{(k)}_{y,x}))v$ of the unit blocks. $P(A)\defeq Q(A)^*Q(A)$, with $Q(A)$ the covariant average over the $L$-blocks of the unit lattice, is the orthogonal projection onto the transported constants of the $L$-blocks, with $|P(A)(y,y')|\le1$ and $P(A)(y,y')=0$ unless $y,y'$ lie in one $L$-block. The operators of this part are
\[
G_k(\Omega,\Lambda,A)\defeq T_\Lambda^{-1},\quad T_\Lambda\defeq-\Delta^{\eta,N}_{A,\Omega}+m_k^2+a_kP_k(A)(\Omega\setminus B^k\Lambda)+a_{k+1}L^{-2}P_{k+1}(A)(B^k\Lambda)
\]
\cite[(3.5)]{Ba4}, for $\Lambda\subset\Omega^{(k)}$ a union of $L$-blocks ($P_k(A)(Y)$ is the sum of the block projections of the unit blocks in $Y$, $P_{k+1}(A)(Y)$ that of the $L$-blocks), so that $T_\emptyset=G_k(\Omega,A)^{-1}$ and $T_{\Omega^{(k)}}$ is the operator of $G^\eta_{k+1}(\Omega,A)$; and
\[
X_\Lambda\defeq(\Delta^{(k)}(\Omega,A)+aL^{-2}P(A))|_\Lambda,\qquad C_\Lambda=X_\Lambda^{-1}
\]
\cite[(1.13)]{Ba4}, with $X|_\Lambda\defeq\Lambda X\Lambda$ on $\ell^2(\Lambda)$. Three facts about Ba{\l}aban's blocks and contours are used. (H1) $\Omega^{(k)}$ and $\Lambda$ are unions of large blocks, and every large block is a union of $L$-blocks (\cite[p.~607]{Ba1}: the large-block partition is the block partition with $L$ replaced by $ML$). (H2) The contours are composed, $\Gamma^{(k+1)}_{z,x}=\Gamma_{z,y}\cup\Gamma^{(k)}_{y,x}$ for $x\in B^k(y)$, $y\in B(z)$ \cite[(2.2)]{Ba1}; hence $Q_{k+1}(A)=Q(A)Q_k(A)$ and $P_{k+1}(A)=Q_k^*P(A)Q_k$ is the orthogonal projection onto the transported constants of the $L$-blocks $B^{k+1}(z)$. (H3) Each contour is a union of axis-ordered monotone lattice paths inside its block \cite[(2.1)]{Ba1}; the $\eta$-length of $\Gamma^{(k)}_{y,x}$ is at most $d(L-1)L^{-1}(1+L^{-1}+\dots)\le d$, that of $\Gamma^{(k+1)}_{z,x}$ at most $d(L-1)+d\le d(L+1)$. A vector field $A$ is \emph{regular} on $\Omega$ when $|\partial^\eta_\mu A_\nu(x)|\le ce_k^{\beta-1}$ for $x\in\Omega$, some $\beta>0$ and a universal $c$ \cite[(1.7)]{Ba4}; its curvature $F_{\nu\mu}\defeq\partial^\eta_\nu A_\mu-\partial^\eta_\mu A_\nu$ \cite[(1.6)]{Ba1} is gauge invariant and then obeys $|F_{\nu\mu}|\le2ce_k^{\beta-1}$. Every smallness condition on $e_k=et_k^{1/2}\le e\eps_0^{1/2}$ below is a cap on $\eps_0$ depending on $(d,N,\norm q,c,\beta,a,L)$ only. For $\theta>0$ put $S(\theta)\defeq\sum_{n\in\Z^d}e^{-\theta|n|}<\infty$.

\begin{lemma}[cube coercivity]\label{lem:a06A:cube}
Let $\Delta\subset\eta\Z^d$ be a cube of side $\ell\in\{1,L\}$ that is a union of unit blocks, $A$ a bond field on $\Delta$ with $|F_{\nu\mu}(A)|\le f$ on the plaquettes of $\Delta$, $\Pi^A$ the orthogonal projection onto $\{x\mapsto U(A(\Gamma_x))v:v\in\R^N\}$ for contours $\Gamma_x\subset\Delta$ from a root, each of $\eta$-length at most $\lambda_\Gamma$, and $c_p>0$. Put $\alpha\defeq(d-1)\ell f$ and assume $\norm qe_k\max\{\lambda_\Gamma,1\}\alpha\le\ln2$. Then, with $\kappa\defeq2\norm qe_k\alpha$ and $\eps_g\defeq2\norm qe_k\lambda_\Gamma\alpha$,
\[
-\Delta^{\eta,N}_{A,\Delta}+c_p\Pi^A\ \ge\ \min\{2\ell^{-2},c_p\}-2c_p\eps_g-d\kappa^2 .
\]
If $A$ is regular on $\Delta$, then $f\le2ce_k^{\beta-1}$ and $\kappa,\eps_g=O(e_k^\beta)$, so there is $e_1(d,N,\norm q,c,\beta,L,a)>0$ such that for $e_k\le e_1$ the right side is at least $\frac12\min\{2\ell^{-2},c_p\}$ for every $c_p\le a$ and every $\lambda_\Gamma\le d(L+1)$.
\end{lemma}

\begin{proof}
\emph{Gauge covariance.} For $\lambda\colon\Delta\to\R$ the gauge transform $A^\lambda_b\defeq A_b+\eta^{-1}(\lambda(b_+)-\lambda(b_-))$ gives $U(A^\lambda_b)=U(A_b)e^{qe_k(\lambda(b_+)-\lambda(b_-))}$. With the orthogonal multiplication operator $G\defeq e^{-qe_k\lambda}$ (acting on bond functions through $\lambda(b_-)$), $D_{A^\lambda}=GD_AG^{-1}$, since all the matrices commute; and $U(A^\lambda(\Gamma_x))=e^{qe_k\lambda(x)}U(A(\Gamma_x))e^{-qe_k\lambda(x_0)}$ for the root $x_0$, so the range of $\Pi^{A^\lambda}$ is $G^{-1}$ of the range of $\Pi^A$ and $\Pi^{A^\lambda}=G\Pi^AG^{-1}$. Hence $-\Delta^{\eta,N}_{A^\lambda,\Delta}+c_p\Pi^{A^\lambda}=G(-\Delta^{\eta,N}_{A,\Delta}+c_p\Pi^A)G^{-1}$, and the lower bound is gauge invariant; $F$ is unchanged.

\emph{Axial gauge.} Choose $\lambda$ so that $A'\defeq A^\lambda$ has $A'_1\equiv0$, $A'_2=0$ on the bonds of one face $\{x_1=x_1^0\}$, $A'_3=0$ on the bonds of one line $\{x_1=x_1^0,x_2=x_2^0\}$, and so on in $d$ dimensions (the bonds on which a component is prescribed form a tree, so $\lambda$ exists). Then $A'_\mu(x)$ is reached from a bond where it vanishes by at most $d-1$ axis-parallel segments of length at most $\ell$, along each of which it changes by $\eta$ times the curvature: on a segment in direction $\nu<\mu$, $\partial^\eta_\nu A'_\mu=F_{\nu\mu}+\partial^\eta_\mu A'_\nu$, and $A'_\nu$ vanishes identically on the set over which the segment runs. Hence $|A'_\mu|\le(d-1)\ell f=\alpha$ on $\Delta$.

\emph{Bond and pin defects.} $|U(A'_b)-1|\le\norm qe_k\eta|A'_b|e^{\norm qe_k\eta|A'_b|}\le2\norm qe_k\alpha\eta$, so $\kappa\ge\sup_b|U(A'_b)-1|/\eta$. The transported constants are $g(x)\defeq U(A'(\Gamma_x))=\exp(qe_k\sum_{b\in\Gamma_x}\eta A'_b)$, and $\sum_{b\in\Gamma_x}\eta|A'_b|\le\lambda_\Gamma\alpha$ gives $|g(x)-1|\le2\norm qe_k\lambda_\Gamma\alpha=\eps_g$.

\emph{Error operator.} Write $D_{A'}=D_1+E$ with $(D_1\phi)(b)=\eta^{-1}(\phi(b_+)-\phi(b_-))$ and $(E\phi)(b)=\eta^{-1}(U(A'_b)-1)\phi(b_+)$. Each site is the final point of at most $d$ bonds of $\Delta$, so $\norm{E\phi}^2\le\kappa^2\sum_b\eta^d|\phi(b_+)|^2\le d\kappa^2\norm\phi^2$, and $\norm{D_{A'}\phi}^2\ge\frac12\norm{D_1\phi}^2-d\kappa^2\norm\phi^2$.

\emph{Neumann gap.} $\Delta$ is a product of $d$ paths with $\ell/\eta$ sites and spacing $\eta$; the Neumann Laplacian of a path has eigenvalues $4\eta^{-2}\sin^2(\pi j\eta/(2\ell))$, $0\le j<\ell/\eta$, so $\norm{D_1\phi}^2\ge\lambda_\ell\norm{(1-\Pi_1)\phi}^2$ with $\Pi_1$ the projection onto the constants and $\lambda_\ell=4\eta^{-2}\sin^2(\pi\eta/(2\ell))\ge4\ell^{-2}$, by $\sin t\ge2t/\pi$ on $[0,\pi/2]$.

\emph{Pin comparison.} $\Pi^{A'}=V_gV_g^*$ with $(V_gv)(x)\defeq g(x)v/\ell^{d/2}$, an isometry $\R^N\to\ell^2(\Delta)$, and $\Pi_1=V_1V_1^*$; $\norm{V_g-V_1}\le\sup_x|g(x)-1|\le\eps_g$, so $\norm{\Pi^{A'}-\Pi_1}\le2\eps_g$.

\emph{Combination.} $-\Delta^{\eta,N}_{A',\Delta}+c_p\Pi^{A'}\ge\frac12\lambda_\ell(1-\Pi_1)+c_p\Pi_1-2c_p\eps_g-d\kappa^2\ge\min\{2\ell^{-2},c_p\}-2c_p\eps_g-d\kappa^2$. The last sentence of the statement follows since $\alpha\le2(d-1)Lce_k^{\beta-1}$ and $\kappa,\eps_g\le4d(L+1)\norm q(d-1)Lc\,e_k^\beta$.
\end{proof}

\begin{lemma}[{coercivity with mixed pins; \cite[(1.8)]{Ba4} on a general region}]\label{lem:a06A:coer}
Assume (H1)--(H3), $A$ regular on $\Omega$ and $e_k\le e_1$. Then for every $\Lambda\subset\Omega^{(k)}$ that is a union of $L$-blocks ($\Lambda=\emptyset$ and $\Lambda=\Omega^{(k)}$ included)
\[
T_\Lambda\ \ge\ \gamma_L\defeq\tfrac14\min\{2L^{-2},a_*L^{-2}\},
\]
and $T_\emptyset=G_k(\Omega,A)^{-1}\ge\gamma_\emptyset\defeq\frac14\min\{2,a_*\}$, which is \cite[(1.8)]{Ba4} on every union $\Omega$ of unit blocks, with $\gamma_0=\gamma_\emptyset$ independent of $L$, $\eta$, $k$, $\Omega$, $A$ and the torus.
\end{lemma}

\begin{proof}
Partition $\Omega$ into the unit blocks of $\Omega\setminus B^k\Lambda$ and the $L$-blocks of $B^k\Lambda$ (cells). Dropping the bonds between cells removes non-negative terms from the Dirichlet form, so $-\Delta^{\eta,N}_{A,\Omega}\ge\bigoplus_{\text{cells}}(-\Delta^{\eta,N}_{A,\text{cell}})$. The pins are block-diagonal on this partition: $P_k(A)(\Omega\setminus B^k\Lambda)$ is the sum of the projections $\Pi^A$ of the unit blocks, with contours of $\eta$-length at most $d$ (H3), and $P_{k+1}(A)(B^k\Lambda)$ is the sum of the projections of the $L$-blocks (H2), with contours of $\eta$-length at most $d(L+1)$ (H3). On each cell apply \cref{lem:a06A:cube} with $(\ell,c_p)=(1,a_k)$, resp. $(L,a_{k+1}L^{-2})$, and $c_p\le a$: the cell operator is at least $\frac12\min\{2,a_*\}\ge\gamma_L$, resp. $\frac12\min\{2L^{-2},a_*L^{-2}\}\ge\gamma_L$. Drop $m_k^2\ge0$. For $\Lambda=\emptyset$ all cells are unit blocks and the bound is $\frac12\min\{2,a_*\}\ge\gamma_\emptyset$.
\end{proof}

\begin{lemma}[{Combes--Thomas bound; the $L^2$ form of \cite[Cor.~2.3]{Ba4} for every pin structure}]\label{lem:a06A:CT}
Let $T\defeq D_A^*D_A+V+\sum_Bc_B\Pi_B$ on $\ell^2(\Omega,\eta^d)$, where $V\ge0$ is a multiplication operator, the $\Pi_B$ are orthogonal projections whose ranges are supported in pairwise disjoint sets $B$ of diameter at most $\ell_*$, $0\le c_B\le a$, and $T\ge\gamma>0$. Put
\[
\theta_0(\gamma,\ell_*)\defeq\min\Big\{\frac1{\ell_*},\ \frac{\gamma^{1/2}}{36\sqrt d},\ \frac{\gamma}{12\mathrm e^2a\ell_*}\Big\},
\]
$\mathrm e$ Euler's number. Then for $f,f'\in\ell^2(\Omega)$ and bond functions $g,g'$, with $\delta$ the distance of the supports (bonds located at $b_-$) and $\theta\defeq\theta_0(\gamma,\ell_*)$:
\begin{enumerate}[label=\textup{(\roman*)},nosep,leftmargin=*]
\item $|\langle f,T^{-1}f'\rangle|\le2\gamma^{-1}e^{-\theta\delta}\norm f\norm{f'}$;
\item $|\langle g,D_AT^{-1}f'\rangle|$ and $|\langle f,T^{-1}D_A^*g'\rangle|$ are at most $(2\gamma^{-1/2}+6\sqrt d\,\theta\gamma^{-1})e^{\theta}e^{-\theta\delta}$ times the product of the norms;
\item $|\langle g,D_AT^{-1}D_A^*g'\rangle|\le4e^{\theta}e^{-\theta\delta}\norm g\norm{g'}$.
\end{enumerate}
The constants do not depend on $\eta$, $\Omega$, $A$ or the torus. $\theta_0(\gamma,\ell_*)$ is increasing in $\gamma$ and decreasing in $\ell_*$; in particular $\theta_0(\gamma_\emptyset,1)\ge\theta_0(\gamma_L,L)$.
\end{lemma}

\begin{proof}
Let $\rho\defeq\dist(\cdot,\supp f')$ (for bond functions evaluated at $b_-$), a $1$-Lipschitz function along bonds, and $T_\theta\defeq e^{\theta\rho}Te^{-\theta\rho}$ for $0<\theta\le\theta_0(\gamma,\ell_*)$; then $T^{-1}=e^{-\theta\rho}T_\theta^{-1}e^{\theta\rho}$ and $e^{\theta\rho}f'=f'$.

\emph{Hopping.} $(e^{\theta\rho}D_Ae^{-\theta\rho}\phi)(b)=(D_A\phi)(b)+\eta^{-1}(e^{-\theta(\rho(b_+)-\rho(b_-))}-1)U(A_b)\phi(b_+)$. Since $|\rho(b_+)-\rho(b_-)|\le\eta$ and $\theta\eta\le1$, $|\eta^{-1}(e^{-\theta s}-1)|\le3\theta$ for $|s|\le\eta$, so $e^{\theta\rho}D_Ae^{-\theta\rho}=D_A+\theta B_1$ with $\norm{B_1}\le3\sqrt d$ (each site is the final point of at most $d$ bonds). Likewise $e^{\theta\rho}D_A^*e^{-\theta\rho}=(e^{-\theta\rho}D_Ae^{\theta\rho})^*=(D_A-\theta B_1')^*$ with $\norm{B_1'}\le3\sqrt d$. Hence $e^{\theta\rho}D_A^*D_Ae^{-\theta\rho}-D_A^*D_A=\theta(D_A^*B_1-B_1'^*D_A)-\theta^2B_1'^*B_1$, and
\[
|\langle\phi,(e^{\theta\rho}D_A^*D_Ae^{-\theta\rho}-D_A^*D_A)\psi\rangle|\le3\sqrt d\,\theta(\norm{D_A\phi}\norm\psi+\norm\phi\norm{D_A\psi})+9d\theta^2\norm\phi\norm\psi .
\]
$V$ commutes with $e^{\theta\rho}$.

\emph{Pins.} Write $\Pi_B=VV^*$ with $V$ an isometry onto the range, and let $\rho_B$ be the value of $\rho$ at a point of $B$. Then $e^{\theta\rho}\Pi_Be^{-\theta\rho}=(e^{\theta(\rho-\rho_B)}V)(e^{-\theta(\rho-\rho_B)}V)^*$, and $|\rho-\rho_B|\le\ell_*$ on $B$, so $\norm{e^{\pm\theta(\rho-\rho_B)}V-V}\le\theta\ell_*e^{\theta\ell_*}$ and $\norm{e^{\theta\rho}\Pi_Be^{-\theta\rho}-\Pi_B}\le2\theta\ell_*e^{\theta\ell_*}+\theta^2\ell_*^2e^{2\theta\ell_*}\le2\theta\ell_*e^{2\theta\ell_*}$, using $te^t\le2(e^t-1)$ for $t=\theta\ell_*\le1$. The sets $B$ are disjoint, so the pin changes are block-diagonal and their sum has norm at most $2a\theta\ell_*e^{2\theta\ell_*}\le2\mathrm e^2a\theta\ell_*$.

\emph{Perturbation.} Put $K\defeq T^{-1/2}(T_\theta-T)T^{-1/2}$. Since $D_A^*D_A\le T$, $\norm{D_AT^{-1/2}}\le1$, and $\norm{T^{-1/2}}\le\gamma^{-1/2}$, so
\[
\norm K\le6\sqrt d\,\theta\gamma^{-1/2}+9d\theta^2\gamma^{-1}+2\mathrm e^2a\theta\ell_*\gamma^{-1}\le\tfrac16+\tfrac1{144}+\tfrac16<\tfrac12
\]
for $\theta\le\theta_0(\gamma,\ell_*)$. Hence $T_\theta=T^{1/2}(1+K)T^{1/2}$ is invertible, $T_\theta^{-1}=T^{-1/2}(1+K)^{-1}T^{-1/2}$, and $\norm{T_\theta^{-1}}\le2\gamma^{-1}$, $\norm{D_AT_\theta^{-1}}\le2\gamma^{-1/2}$, $\norm{D_AT_\theta^{-1}D_A^*}\le2$.

\emph{Bounds.} (i) $\langle f,T^{-1}f'\rangle=\langle e^{-\theta\rho}f,T_\theta^{-1}f'\rangle$ and $\norm{e^{-\theta\rho}f}\le e^{-\theta\delta}\norm f$. (ii) $\langle g,D_AT^{-1}f'\rangle=\langle e^{-\theta\rho}g,(D_A+\theta B_1)T_\theta^{-1}f'\rangle$, bounded by $(2\gamma^{-1/2}+3\sqrt d\,\theta\cdot2\gamma^{-1})e^{-\theta\delta}\norm g\norm{f'}$; the adjoint form follows from the symmetry of $T$ with $\rho$ the distance to $\supp f$. (iii) With $\rho$ the distance to $\supp g'$, $e^{\theta\rho}g'=g'$ and $\langle g,D_AT^{-1}D_A^*g'\rangle=\langle e^{-\theta\rho}g,(D_A+\theta B_1)T_\theta^{-1}(D_A-\theta B_1')^*g'\rangle$, bounded by $(2+12\sqrt d\,\theta\gamma^{-1/2}+18d\theta^2\gamma^{-1})e^{-\theta\delta}\norm g\norm{g'}\le4e^{-\theta\delta}\norm g\norm{g'}$. The factors $e^\theta$ in (ii)--(iii) are $\ge1$. The monotonicity of $\theta_0$ is read off its definition, and $\gamma_\emptyset\ge\gamma_L$, $1\le L$.
\end{proof}

With \cref{lem:a06A:coer} ($\ell_*=L$, $\gamma=\gamma_L$: the pins of $T_\Lambda$ are the unit-block and $L$-block projections, of diameters $\le1$ and $\le L$, with $c_B\in\{a_k,a_{k+1}L^{-2}\}\subset[0,a]$), \cref{lem:a06A:CT} gives \cite[Cor.~2.3]{Ba4} for $G_k(\Omega,\Lambda,A)$ and $G^\eta_{k+1}(\Omega,A)$, with constants depending on $L$ as \cite[p.~588]{Ba4} says; with $\Lambda=\emptyset$ ($\ell_*=1$, $\gamma=\gamma_\emptyset$) it is a proof of the first clause of \cite[Cor.~2.3]{Ba4}, the bounds for $G_k(\Omega,A)$ (its $\delta G_k$ clause is not used in this appendix). Throughout Steps~3(B)--3(D), $\theta_0\defeq\theta_0(\gamma_L,L)$; by the last sentence of \cref{lem:a06A:CT} every Combes--Thomas rate of this part is at least $\theta_0$.

% a06_BALtorus_A2: continuation of a06_BALtorus_A1 (input from it). Group (A), part 2 of 4:
% the block inversion (3.8) as an identity, (1.15), (1.16); the resolvent difference (3.9)-(3.10); the conditioning-region difference (1.18).

\begin{lemma}[{block inversion; \cite[(3.8), (1.15), (1.16)]{Ba4}}]\label{lem:a06A:blockinv}
Assume (H1)--(H3), $A$ regular on $\Omega$, $e_k\le e_1$, and let $\Lambda\subset\Omega^{(k)}$ be a union of $L$-blocks. Put $P_\Lambda\defeq\Lambda P(A)\Lambda$ (an orthogonal projection on $\ell^2(\Lambda)$, since $P(A)$ is block-diagonal over the $L$-blocks) and $W\defeq Q_k^*\Lambda(1-(a_{k+1}/a_k)L^{-2}P_\Lambda)$. Then:
\begin{enumerate}[label=\textup{(\alph*)},nosep,leftmargin=*]
\item $X_\Lambda$ is invertible and $X_\Lambda^{-1}=W^*G_k(\Omega,\Lambda,A)W+a_k^{-1}-(a_{k+1}L^{-2}/a_k^2)P_\Lambda$; that is, $W\delta_y=q_k(y)-(a_{k+1}/a_k)L^{-2}q_{k+1}(y)$ with $q_{k+1}(y)\defeq Q_k^*P(A)\delta_y$, which is \cite[(3.8)]{Ba4}.
\item $\gamma_X\le X_\Lambda\le a_k+aL^{-2}\le\gamma_1\defeq a+aL^{-2}$, with $\gamma_X\defeq(a_*^{-1}+\gamma_L^{-1})^{-1}$: this is \cite[(1.15)]{Ba4}.
\item $|C_\Lambda(y,y')|\le c_2e^{-\theta_0|y-y'|}$ for $y,y'\in\Lambda$, with $c_2\defeq(2\gamma_L^{-1}+2a_*^{-1})e^{2\theta_0L}$ and $\theta_0=\theta_0(\gamma_L,L)$: this is \cite[(1.16)]{Ba4}, also for $C^{(k)}(\Omega,A)=C_{\Omega^{(k)}}$.
\end{enumerate}
The constants depend on $(d,N,\norm q,a,L)$ only.
\end{lemma}

\begin{proof}
On $\ell^2(\Lambda)\oplus\ell^2(\Omega)$ let $\mathcal M\defeq\begin{pmatrix}A_\psi&-B\\-B^*&D_\phi\end{pmatrix}$ with $A_\psi\defeq a_k+aL^{-2}P_\Lambda$, $B\defeq a_k\Lambda Q_k$, $D_\phi\defeq H+a_kP_k(A)$, $H\defeq-\Delta^{\eta,N}_{A,\Omega}+m_k^2$. Its quadratic form is
\[
a_k\norm{\psi-\Lambda Q_k\phi}^2+a_k\norm{(1-\Lambda)Q_k\phi}^2+aL^{-2}\norm{P_\Lambda\psi}^2+\langle\phi,H\phi\rangle .
\]
\emph{Schur complement in $\phi$.} $D_\phi=G_k(\Omega,A)^{-1}\ge\gamma_\emptyset$ by \cref{lem:a06A:coer}, and $A_\psi-BD_\phi^{-1}B^*=\Lambda(a_k-a_k^2Q_kG_k(\Omega,A)Q_k^*+aL^{-2}P(A))\Lambda=X_\Lambda$ by \cite[(1.14)]{Ba4}.

\emph{Schur complement in $\psi$.} $A_\psi\ge a_k>0$, and with $b\defeq aL^{-2}$, $(a_k+bP_\Lambda)^{-1}=a_k^{-1}(1-P_\Lambda)+(a_k+b)^{-1}P_\Lambda$, so $a_k-a_k^2(a_k+bP_\Lambda)^{-1}=\frac{a_kb}{a_k+b}P_\Lambda=a_{k+1}L^{-2}P_\Lambda$ by the recursion $a_{k+1}=aa_k/(aL^{-2}+a_k)$. Hence
\[
D_\phi-B^*A_\psi^{-1}B=H+a_kQ_k^*(1-\Lambda)Q_k+a_{k+1}L^{-2}Q_k^*P_\Lambda Q_k=T_\Lambda,
\]
since $Q_k^*(1-\Lambda)Q_k=P_k(A)(\Omega\setminus B^k\Lambda)$ and $Q_k^*P_\Lambda Q_k=P_{k+1}(A)(B^k\Lambda)$ by (H2). This is the computation \cite[(3.4)--(3.5)]{Ba4} done as algebra.

\emph{Block inverse.} $A_\psi>0$ and $T_\Lambda\ge\gamma_L>0$ (\cref{lem:a06A:coer}), so $\mathcal M$ is positive definite, and its inverse's $\psi\psi$ block is both $(A_\psi-BD_\phi^{-1}B^*)^{-1}=X_\Lambda^{-1}$ and $A_\psi^{-1}+A_\psi^{-1}BT_\Lambda^{-1}B^*A_\psi^{-1}$. Now $A_\psi^{-1}B=(1-\frac{b}{a_k+b}P_\Lambda)\Lambda Q_k=W^*$, because $\frac{b}{a_k+b}=\frac{a_{k+1}}{a_k}L^{-2}$, and $A_\psi^{-1}=a_k^{-1}-\frac{b}{a_k(a_k+b)}P_\Lambda=a_k^{-1}-\frac{a_{k+1}L^{-2}}{a_k^2}P_\Lambda$. This is (a); $P(A)\delta_y$ is supported in the $L$-block of $y$, which lies in $\Lambda$ by (H1), so $\Lambda P_\Lambda\delta_y=P(A)\delta_y$ and $W\delta_y$ is as stated.

\emph{(b).} $\norm{Q_k}\le1$ (Cauchy--Schwarz on blocks of volume $1$) and $\norm{1-tP_\Lambda}\le1$ for $t=(a_{k+1}/a_k)L^{-2}\in[0,1]$ (as $a_{k+1}\le a_kL^2$), so $\norm W\le1$ and $\norm{X_\Lambda^{-1}}\le a_k^{-1}+\gamma_L^{-1}\le a_*^{-1}+\gamma_L^{-1}$. The upper bound holds because $a_k^2Q_kG_kQ_k^*\ge0$ and $\norm{P_\Lambda}\le1$.

\emph{(c).} $W\delta_y$ is supported in the $L$-block $B^{k+1}$ containing $y$, a cube of side $L$ containing $y$, and $\norm{W\delta_y}\le1$. \Cref{lem:a06A:CT}(i) for $T_\Lambda$ (pins of diameter $\le L$, $\gamma=\gamma_L$) gives $|\langle W\delta_y,T_\Lambda^{-1}W\delta_{y'}\rangle|\le2\gamma_L^{-1}e^{-\theta_0(|y-y'|-2L)}$. The term $a_k^{-1}\delta_{yy'}$ is at most $a_*^{-1}$ and vanishes unless $y=y'$; the term $(a_{k+1}L^{-2}/a_k^2)P_\Lambda(y,y')$ has $a_{k+1}L^{-2}/a_k^2=\frac{b}{a_k(a_k+b)}\le a_*^{-1}$, $|P_\Lambda(y,y')|\le1$, and vanishes unless $|y-y'|\le L$. Summing, $|C_\Lambda(y,y')|\le(2\gamma_L^{-1}e^{2\theta_0L}+a_*^{-1}+a_*^{-1}e^{\theta_0L})e^{-\theta_0|y-y'|}\le c_2e^{-\theta_0|y-y'|}$. $\Lambda=\Omega^{(k)}$ is allowed.
\end{proof}

\begin{lemma}[{the resolvent difference; \cite[(3.9), (3.10)]{Ba4}}]\label{lem:a06A:resdiff}
Under the hypotheses of \cref{lem:a06A:blockinv}:
\begin{enumerate}[label=\textup{(\alph*)},nosep,leftmargin=*]
\item $G_k(\Omega,\Lambda,A)-G^\eta_{k+1}(\Omega,A)=G_k(\Omega,\Lambda,A)\,M'\,G^\eta_{k+1}(\Omega,A)$ with $M'\defeq[a_{k+1}L^{-2}P_{k+1}(A)-a_kP_k(A)](\Omega\setminus B^k\Lambda)$.
\item For $f,f'\in\ell^2(B^k\Lambda)$, with $d_f\defeq\dist(\supp f,\Omega\setminus B^k\Lambda)$,
\[
|\langle f,(G_k(\Omega,\Lambda,A)-G^\eta_{k+1}(\Omega,A))f'\rangle|\le c_3\norm f\norm{f'}\exp\big[-\tfrac{\theta_0}6\big(\dist(\supp f,\supp f')+d_f+d_{f'}\big)\big],
\]
with $c_3\defeq8a\gamma_L^{-2}e^{\theta_0L/3}S_L(\theta_0/3)S_L(\theta_0/6)$, $S_L(\theta)\defeq\sum_{n\in\Z^d}e^{-\theta L(|n|-1)_+}$. This is stronger than \cite[(3.10)]{Ba4}, whose $B^k(\Lambda^c)$ contains $\Omega\setminus B^k\Lambda$.
\end{enumerate}
\end{lemma}

\begin{proof}
(a) is the second resolvent identity, since $T_{\Omega^{(k)}}-T_\Lambda=M'$. (The bracket of \cite[(3.9)]{Ba4} is printed with the opposite sign; only norms of it are used there.)

(b) $M'$ is block-diagonal over the $L$-blocks $\Delta\subset\Omega\setminus B^k\Lambda$ (by (H2) each $P_{k+1}$ summand and the $P_k$ summands it contains live on one $L$-block), with $\norm{M'_\Delta}\le a_{k+1}L^{-2}+a_k\le2a$. Let first $f$, $f'$ be supported in single $L$-blocks $B,B'\subset B^k\Lambda$, and write $u_\Delta\defeq\dist(\Delta,\supp f)$, $v_\Delta\defeq\dist(\Delta,\supp f')$. Both $T_\Lambda$ and $T_{\Omega^{(k)}}$ are at least $\gamma_L$ (\cref{lem:a06A:coer}), so \cref{lem:a06A:CT}(i) gives $\norm{\mathbf 1_\Delta G_k(\Omega,\Lambda,A)f}\le2\gamma_L^{-1}e^{-\theta_0u_\Delta}\norm f$ and the same for $G^\eta_{k+1}(\Omega,A)f'$ with $v_\Delta$; hence
\[
|\langle f,(G_k(\Omega,\Lambda,A)-G^\eta_{k+1}(\Omega,A))f'\rangle|\le2a(2\gamma_L^{-1})^2\sum_\Delta e^{-\theta_0(u_\Delta+v_\Delta)}\norm f\norm{f'} .
\]
Since $\Delta\subset\Omega\setminus B^k\Lambda$, $u_\Delta\ge d_f$ and $v_\Delta\ge d_{f'}$; since $\Delta$ has diameter $L$, $\dist(\supp f,\supp f')\le u_\Delta+L+v_\Delta$. Therefore $u_\Delta+v_\Delta\ge\frac13(\dist(\supp f,\supp f')-L)+\frac13(d_f+d_{f'})+\frac13(u_\Delta+v_\Delta)$, and $\sum_\Delta e^{-(\theta_0/3)u_\Delta}\le S_L(\theta_0/3)$, a sum over the $L$-block lattice, uniform in the torus (on a torus, each block has a nearest lift, and distinct blocks have distinct lifts). So for block-supported $f,f'$ the bound holds with the rate $\theta_0/3$ and the constant $8a\gamma_L^{-2}e^{\theta_0L/3}S_L(\theta_0/3)$. For general $f,f'$ decompose $f=\sum_Bf_B$, $f'=\sum_{B'}f'_{B'}$ over the $L$-blocks of $B^k\Lambda$; restricting supports increases all three distances, so $\dist(\supp f_B,\supp f'_{B'})+d_{f_B}+d_{f'_{B'}}\ge\dist(\supp f,\supp f')+d_f+d_{f'}$, and splitting the rate $\theta_0/3$ into two halves, one half gives the exponential of the statement and the other leaves the kernel $e^{-(\theta_0/6)\dist(\supp f_B,\supp f'_{B'})}\le e^{-(\theta_0/6)\dist(B,B')}$, whose Schur norm on $\ell^2$ of the block indices is at most $S_L(\theta_0/6)$; since $\sum_B\norm{f_B}^2=\norm f^2$, the Schur test gives (b).
\end{proof}

\begin{lemma}[{the conditioning-region difference; \cite[(1.18)]{Ba4}}]\label{lem:a06A:CLdiff}
Under the hypotheses of \cref{lem:a06A:blockinv}, for $y,y'\in\Lambda$,
\[
|C_\Lambda(y,y')-C^{(k)}(\Omega,A;y,y')|\le c_4\exp\big[-\tfrac{\theta_0}6\big(|y-y'|+\dist(y,\Omega^{(k)}\setminus\Lambda)+\dist(y',\Omega^{(k)}\setminus\Lambda)\big)\big],
\]
with $c_4\defeq c_3e^{(\theta_0/6)(4L+2)}$. The bound decays a fortiori in the distances to the whole complement $\Lambda^c$.
\end{lemma}

\begin{proof}
$C^{(k)}(\Omega,A)=C_{\Omega^{(k)}}$. For $y\in\Lambda$, $W_{\Omega^{(k)}}\delta_y=W_\Lambda\delta_y=Q_k^*(\delta_y-(a_{k+1}/a_k)L^{-2}P(A)\delta_y)$, because $P(A)\delta_y$ lives on the $L$-block of $y$, which lies in $\Lambda$ (H1); and $P_\Lambda(y,y')=P(A)(y,y')$ for $y,y'\in\Lambda$. By \cref{lem:a06A:blockinv}(a) for both $\Lambda$ and $\Omega^{(k)}$ the constant terms cancel and
\[
C_\Lambda(y,y')-C_{\Omega^{(k)}}(y,y')=\langle W\delta_y,(G_k(\Omega,\Lambda,A)-G^\eta_{k+1}(\Omega,A))W\delta_{y'}\rangle .
\]
Apply \cref{lem:a06A:resdiff}(b) with $f=W\delta_y$, $f'=W\delta_{y'}$, both of norm $\le1$ and supported in $B^k$ of the $L$-blocks of $y$, $y'$, which lie in $B^k\Lambda$. Every point of $\supp f$ is within $L$ of $y$, and $\Omega\setminus B^k\Lambda=B^k(\Omega^{(k)}\setminus\Lambda)$ lies within $1$ of $\Omega^{(k)}\setminus\Lambda$; so $\dist(\supp f,\supp f')\ge|y-y'|-2L$ and $d_f\ge\dist(y,\Omega^{(k)}\setminus\Lambda)-L-1$, and the shifts total at most $4L+2$.
\end{proof}

% a06_BALtorus_A3: continuation of a06_BALtorus_A2 (input from it). Group (A), part 3 of 4:
% the two-sided depth factor of Ba4's Theorem, the domain differences (5.5) and (1.20), Ba1 Prop. 2.2, the summary theorem, the polylog bound.

The remaining estimates use the printed theorem of \cite{Ba4}, which is \cite[Prop.~2.1]{Ba1}: for $\alpha<1$ there are $\delta_0,c_0,R_0>0$, independent of $A$, $k$, $\Omega$, depending on $d$, $M$ (and $c_0$ on $\alpha$), such that for $e_k$ small and $f\colon\Omega\to\R^N$: the $C^{1,\alpha}$ bound \cite[(1.9)]{Ba4} for $x,x'\in\Omega$ with $\dist(\{x,x'\},\Omega^c)\ge R_0$; $|(D^\eta_{A,\mu}G_k(\Omega,A)f)(x)|,\,|(G_k(\Omega,A)f)(x)|\le c_0e^{-\delta_0\dist(x,\supp f)}\norm f_\infty$ for $x\in\Omega$ with $\dist(x,\Omega^c)\ge R_0$ \cite[(1.10)]{Ba4}; and, for $\Omega\subset\Omega_0$ and $\delta G_k(\Omega,\Omega_0,A)\defeq G_k(\Omega,A)-G_k(\Omega_0,A)$ \cite[(1.11)]{Ba4}, the same bounds with the same restrictions and an additional factor \cite[(1.12)]{Ba4}, printed in \cite[(2.26)]{Ba1} as $\exp(-\delta_0\dist(\supp f,\Omega^c)-\delta_0\dist(\{x,x'\},\Omega^c))$. When $\Omega$ is the whole torus the depth condition is void \cite[p.~611]{Ba1}. Only this restricted form (points at depth at least $R_0$, or $\Omega$ the whole torus) is used anywhere in this appendix; the clause of \cite[p.~611]{Ba1} and \cite[p.~573]{Ba4} that drops the depth condition on rectangular parallelepipeds is not used.

\begin{lemma}[the depth factor, two-sided]\label{lem:a06A:twosided}
Let $\Omega\subset\Omega_0$ be unions of large blocks, $A$ regular on $\Omega_0$, $e_k$ small, $\supp f$ of diameter at most $s$, and $x\in\Omega$ with $\dist(x,\Omega^c)\ge R_0$. Then
\[
|(\delta G_k(\Omega,\Omega_0,A)f)(x)|\le c_0e^{\delta_0s/2}\exp\big[-\tfrac{\delta_0}2\big(\dist(x,\supp f)+\dist(x,\Omega^c)+\dist(\supp f,\Omega^c)\big)\big]\norm f_\infty,
\]
and the same for $D^\eta_{A,\mu}\delta G_k(\Omega,\Omega_0,A)f$, and for the H\"older quotient of \cite[(1.9)]{Ba4} with $|x-x'|\le1$ and $\dist(\{x,x'\},\Omega^c)\ge R_0$, with $\dist(\{x,x'\},\cdot)$ in place of $\dist(x,\cdot)$. For general $f$ (no bound on the diameter) the bounds hold at the rate $\delta_0/4$ with the constant $c_0e^{3\delta_0/4}S(\delta_0/4)$ ($c_0e^{\delta_0}S(\delta_0/4)$ for the H\"older quotient).
\end{lemma}

\begin{proof}
$\dist(x,\Omega^c)\le\dist(x,\supp f)+s+\dist(\supp f,\Omega^c)$, so $\dist(x,\supp f)+\dist(x,\Omega^c)+\dist(\supp f,\Omega^c)\le2\dist(x,\supp f)+2\dist(\supp f,\Omega^c)+s$. The printed factor, read with both of its terms as $\dist(\supp f,\Omega^c)$, gives $c_0e^{-\delta_0\dist(x,\supp f)}e^{-2\delta_0\dist(\supp f,\Omega^c)}\norm f_\infty$, which is at most the right side; read as printed in \cite[(2.26)]{Ba1}, it gives the right side without the factor $e^{\delta_0s/2}$ and at the rate $\delta_0$. So the bound holds under either reading. For the H\"older quotient the printed factor carries $\dist(\{x,x'\},\cdot)$ and the same arithmetic applies, because $\dist(\{x,x'\},\Omega^c)\le\dist(\{x,x'\},\supp f)+s+\dist(\supp f,\Omega^c)$ (the inequality for $x$ and for $x'$ separately, then the minimum over the two points). For general $f$ decompose $\supp f$ into the unit blocks $\square$ it meets, apply the bound to $f\mathbf 1_\square$ ($s\le1$; $\dist(x,\square\cap\supp f)\ge\dist(x,\supp f)$ and $\dist(\square\cap\supp f,\Omega^c)\ge\dist(\supp f,\Omega^c)$), and sum: with $\square_n$ the unit block of $n\in\Z^d$ and $x\in\square_{n_x}$, $\dist(x,\square_n)\ge(|n-n_x|-1)_+$, so $\sum_\square e^{-(\delta_0/2)\dist(x,\square\cap\supp f)}\le e^{-(\delta_0/4)\dist(x,\supp f)}\sum_{n\in\Z^d}e^{-(\delta_0/4)(|n-n_x|-1)_+}\le e^{-(\delta_0/4)\dist(x,\supp f)}e^{\delta_0/4}S(\delta_0/4)$; for the H\"older quotient $\dist(\{x,x'\},\square_n)\ge\dist(x,\square_n)-1$ costs one more factor $e^{\delta_0/4}$.
\end{proof}

\begin{lemma}[{domain differences; \cite[(5.5), (1.20)]{Ba4}}]\label{lem:a06A:domdiff}
Let $\Omega\subset\Omega_0$ be unions of large blocks, (H1)--(H3), $A$ regular on $\Omega_0$, $e_k\le e_1$ and small enough for the printed theorem. Put
\[
\delta_7\defeq\min\{\delta_0/2,\,\theta_0(\gamma_\emptyset,1)\},\qquad\delta^S_7\defeq\min\{\theta_0,\,\delta_7\}=\min\{\theta_0(\gamma_L,L),\,\delta_0/2\}.
\]
\begin{enumerate}[label=\textup{(\alph*)},nosep,leftmargin=*]
\item For $y,y'\in\Omega^{(k)}$, with $d(y)\defeq\dist(B^k(y),\Omega^c)$,
\[
|(\Delta^{(k)}(\Omega,A)-\Delta^{(k)}(\Omega_0,A))(y,y')|\le c_6\exp\big[-\delta_7\big(|y-y'|+d(y)+d(y')\big)\big],
\]
$c_6\defeq\max\{a^2c_0e^{\delta_0/2},\,4a^2\gamma_\emptyset^{-1}e^{2\theta_0(\gamma_\emptyset,1)R_0}\}e^{\delta_7}$: this is \cite[(5.5)]{Ba4}.
\item For every $\Lambda\subset\Omega^{(k)}$ that is a union of $L$-blocks and $x,x'\in\Lambda$, with $\tilde d(x)\defeq\dist(x,\Omega^{(k)c})$,
\[
|C_\Lambda(\Omega,A;x,x')-C_\Lambda(\Omega_0,A;x,x')|\le c_{20}\exp\big[-\tfrac{\delta^S_7}4\big(|x-x'|+\tilde d(x)+\tilde d(x')\big)\big],
\]
$c_{20}\defeq c_2^2c_6e^{\delta^S_7/2}S(\delta^S_7/2)^2$: this is \cite[(1.20)]{Ba4}.
\end{enumerate}
\end{lemma}

\begin{proof}
(a) $Q_k(A)$ on $\Omega$ and on $\Omega_0$ agree on the blocks $B^k(y)$, $y\in\Omega^{(k)}$, so by \cite[(1.14)]{Ba4} the kernel is $\delta\Delta(y,y')=-a_k^2\langle q_k(y),\delta G_k(\Omega,\Omega_0,A)q_k(y')\rangle$, with $q_k(y),q_k(y')$ supported in $\Omega$. Three cases. If $d(y)\ge R_0$, every $x\in B^k(y)$ has $\dist(x,\Omega^c)\ge R_0$, and \cref{lem:a06A:twosided} with $f=q_k(y')$ ($\norm f_\infty=1$, $s\le1$) gives $|(\delta G_kq_k(y'))(x)|\le c_0e^{\delta_0/2}e^{-(\delta_0/2)(\dist(B^k(y),B^k(y'))+d(y)+d(y'))}$ for $x\in B^k(y)$; pairing with $q_k(y)$, whose entries have modulus $1$ on a block of volume $1$, gives $|\delta\Delta(y,y')|\le a_k^2c_0e^{\delta_0/2}e^{-(\delta_0/2)(\dist(B^k(y),B^k(y'))+d(y)+d(y'))}$. If $d(y')\ge R_0$, use $\delta\Delta(y,y')=\delta\Delta(y',y)^{\mathsf T}$. If $d(y),d(y')<R_0$, apply \cref{lem:a06A:CT}(i) to $G_k(\Omega,A)$ and to $G_k(\Omega_0,A)$ separately (\cref{lem:a06A:coer} with $\Lambda=\emptyset$ on $\Omega$ and on $\Omega_0$; $\ell_*=1$, $\gamma=\gamma_\emptyset$): $|\delta\Delta(y,y')|\le4a_k^2\gamma_\emptyset^{-1}e^{-\theta_0(\gamma_\emptyset,1)\dist(B^k(y),B^k(y'))}\le4a_k^2\gamma_\emptyset^{-1}e^{2\theta_0(\gamma_\emptyset,1)R_0}e^{-\theta_0(\gamma_\emptyset,1)(\dist(B^k(y),B^k(y'))+d(y)+d(y'))}$. Finally $\dist(B^k(y),B^k(y'))\ge|y-y'|-1$.

(b) The inverses are bounded by \cref{lem:a06A:blockinv}(b) on both domains, so $C_\Lambda(\Omega,A)-C_\Lambda(\Omega_0,A)=X_\Lambda(\Omega,A)^{-1}[X_\Lambda(\Omega_0,A)-X_\Lambda(\Omega,A)]X_\Lambda(\Omega_0,A)^{-1}$, and $X_\Lambda(\Omega_0,A)-X_\Lambda(\Omega,A)=-\delta\Delta|_\Lambda$, because $\Lambda\subset\Omega^{(k)}$ and $P(A)$ is the same on $\Lambda$ for both domains. With \cref{lem:a06A:blockinv}(c) for both domains and (a),
\[
|\delta C(x,x')|\le c_2^2c_6\sum_{y,y'\in\Lambda}e^{-\delta^S_7\Sigma},\qquad\Sigma\defeq|x-y|+|y-y'|+d(y)+d(y')+|y'-x'| .
\]
Since $\Omega^c=B^k(\Omega^{(k)c})$ and $\dist(B^k(y),B^k(z))\ge|y-z|-1$, $\dist(y,\Omega^{(k)c})\le d(y)+1$; so $\tilde d(x)\le|x-y|+d(y)+1$, $\tilde d(x')\le|x'-y'|+d(y')+1$, and $|x-x'|+\tilde d(x)+\tilde d(x')\le2\Sigma+2$. Hence $e^{-\delta^S_7\Sigma}\le e^{\delta^S_7/2}e^{-(\delta^S_7/4)(|x-x'|+\tilde d(x)+\tilde d(x'))}e^{-(\delta^S_7/2)(|x-y|+|y'-x'|)}$, and the sums over $y$ and $y'$ give $S(\delta^S_7/2)^2$.
\end{proof}

\begin{proposition}[{\cite[Prop.~2.2, (2.27)--(2.29)]{Ba1}}]\label{prop:a06A:Ba1P22}
Assume (H1)--(H3), $A$ regular on $\Omega_0\supset\Omega$, $e_k$ as in \cref{lem:a06A:domdiff}. Then for $y,y'\in\Omega^{(k)}$
\[
|\Delta^{(k)}(\Omega,A;y,y')|\le c_7e^{-\delta_7|y-y'|},\qquad|\Delta^{(k)}(\Omega,A;y,y')-\Delta^{(k)}(\Omega_0,A;y,y')|\le c_7\exp\big[-\delta_7\big(|y-y'|+\dist(y,\Omega^{(k)c})+\dist(y',\Omega^{(k)c})\big)\big],
\]
with $c_7\defeq\max\{a+2a^2\gamma_\emptyset^{-1}e^{\theta_0(\gamma_\emptyset,1)},\,c_6e^{2\delta_7}\}$. The first bound holds on every union $\Omega$ of unit blocks and uses only \cref{lem:a06A:coer,lem:a06A:CT}; its rate is in fact $\theta_0(\gamma_\emptyset,1)\ge\delta_7$.
\end{proposition}

\begin{proof}
$\Delta^{(k)}(y,y')=a_k\delta_{yy'}-a_k^2\langle q_k(y),G_k(\Omega,A)q_k(y')\rangle$ \cite[(1.14)]{Ba4}. \Cref{lem:a06A:coer} with $\Lambda=\emptyset$ and \cref{lem:a06A:CT}(i) with $\ell_*=1$ give $|\langle q_k(y),G_kq_k(y')\rangle|\le2\gamma_\emptyset^{-1}e^{-\theta_0(\gamma_\emptyset,1)\dist(B^k(y),B^k(y'))}$, and $\dist(B^k(y),B^k(y'))\ge|y-y'|-1$; this is the first bound. The second is \cref{lem:a06A:domdiff}(a) with $d(y)\ge\dist(y,\Omega^{(k)c})-1$.
\end{proof}

\begin{theorem}[the propagator estimates of group (A)]\label{thm:a06A:groupA}
Let $d=3$, $L=5$, (H1)--(H3), $\Omega\subset\Omega_0$ unions of large blocks of $T^{(k)}_1$ (the whole torus included), $\Lambda\subset\Omega^{(k)}$ a union of $L$-blocks, $A$ regular on $\Omega_0$, and $\eps_0$ below the caps of \cref{lem:a06A:cube,lem:a06A:domdiff}. Then \cite[(1.8)]{Ba4} on $\Omega$ (\cref{lem:a06A:coer}); the $G_k$ clause of \cite[Cor.~2.3]{Ba4} for $G_k(\Omega,A)$, $G_k(\Omega,\Lambda,A)$ and $G^\eta_{k+1}(\Omega,A)$ (\cref{lem:a06A:CT}); \cite[Prop.~2.3]{Ba1}, i.e.\ \cite[(1.15), (1.16), (1.18), (1.20)]{Ba4} (\cref{lem:a06A:blockinv}(b), (c), \cref{lem:a06A:CLdiff}, \cref{lem:a06A:domdiff}(b)); \cite[(5.5)]{Ba4} (\cref{lem:a06A:domdiff}(a)); and \cite[Prop.~2.2]{Ba1} (\cref{prop:a06A:Ba1P22}) hold with the constants
\[
\gamma_\emptyset,\ \gamma_L,\ \gamma_X,\ \gamma_1,\ \theta_0=\theta_0(\gamma_L,L),\ \theta_0(\gamma_\emptyset,1),\ c_2,\ c_3,\ c_4,\ c_6,\ c_7,\ c_{20},\ \delta_7,\ \delta^S_7
\]
defined above. They depend on $(d,N,\norm q,c,\beta,a,L)$, on $M$ through the constants $\delta_0,c_0,R_0$ of \cite[Thm.]{Ba4} (which enter \cref{lem:a06A:twosided,lem:a06A:domdiff}, \cref{prop:a06A:Ba1P22} and nothing else), and on the cap on $\eps_0$; they do not depend on $\eta$, $k$, $\Omega$, $\Omega_0$, $\Lambda$, $A$, $\eps$ or the torus.
\end{theorem}

\begin{proof}
Each estimate is the lemma named after it; the dependences are read off the definitions of the constants. The order of the arguments is not circular: \cref{lem:a06A:cube,lem:a06A:coer,lem:a06A:CT,lem:a06A:blockinv,lem:a06A:resdiff,lem:a06A:CLdiff} and the first bound of \cref{prop:a06A:Ba1P22} use no estimate of \cite{Ba4} beyond its definitions, while \cref{lem:a06A:twosided,lem:a06A:domdiff} and the second bound of \cref{prop:a06A:Ba1P22} use the printed theorem of \cite{Ba4} in its restricted form (supplied on a torus, with the constants of \cite{Ba4}, by \cref{lem:a06A:torus} below), whose proof \cite[Sect.~2]{Ba4} rests on \cite[Lemmas~2.1, 2.2, 2.4]{Ba4}; the $j$-uniform form of \cite[Lemma~2.4]{Ba4} is \cref{thm:a06:A5}, whose clause \cite[(2.37)]{Ba4} is \cref{lem:a06A:blockinv}(c) at $A=0$, $\Lambda=\Omega^{(k)}$.
\end{proof}

\begin{remark}[dependence on $L$; the depth margin]\label{rem:a06A:L}
(i) $\gamma_\emptyset=\frac14\min\{2,a_*\}$ and the constants of \cref{thm:a06:A5} do not depend on $L$, with $a_*\ge3a/4$ for $L\ge2$; $\gamma_L$, $\gamma_X$, $\theta_0$ and $c_2$ do, through $\gamma_L\propto L^{-2}$, as \cite[p.~588]{Ba4} says of the constants of \cite[(2.37)]{Ba4} and \cite[Prop.~2.3]{Ba1}. At the fixed value $L=5$ all are numbers. (ii) The depth condition of the printed theorem is met wherever it is invoked in Steps~3(B)--3(D): with $r(t)=R(1+\ln t^{-1})^\varrho$ and $R>R_0$ \cite[(2.7)]{Ba2}, every point within distance $1$ of a set at distance more than $2r(t)$ from a boundary has depth more than $2R-1\ge R_0$ (as $R_0\ge1$ may be assumed, enlarging $R_0$ weakening nothing in the theorem); pairs $x,x'$ with $|x-x'|>1$ in \cite[(1.9)]{Ba4} reduce to the sup bounds \cite[p.~578]{Ba4}.
\end{remark}

\begin{lemma}[the polylogarithmic layer]\label{lem:a06A:polylog}
Let $r(t)=R(1+\ln t^{-1})^\varrho$ with $\varrho>1$, and $\delta,R,\kappa>0$. Then $e^{-\delta r(t)}\le t^\kappa$ for all $t\in(0,1]$ if $\delta R\ge\kappa$, and in general for $t\le t_0\defeq\exp(-(\kappa/(\delta R))^{1/(\varrho-1)})$. Consequently $e^{-\delta r(t)}\le C_\kappa t^\kappa$ on $(0,1]$ with $C_\kappa\defeq t_0^{-\kappa}=\exp(\kappa(\kappa/(\delta R))^{1/(\varrho-1)})$.
\end{lemma}

\begin{proof}
With $u\defeq\ln t^{-1}\ge0$, $\delta R(1+u)^\varrho\ge\delta R\max\{1+u,u^\varrho\}$. This is at least $\kappa u$ if $\delta R\ge\kappa$, or if $u^{\varrho-1}\ge\kappa/(\delta R)$, i.e.\ $t\le t_0$. For $t\in[t_0,1]$, $e^{-\delta r(t)}\le1\le t_0^{-\kappa}t^\kappa$.
\end{proof}

% a06_BALtorus_A4: continuation of a06_BALtorus_A3 (input from it). Group (A), part 4 of 4: the torus.

\begin{lemma}[{\cite[Prop.~2.1]{Ba1} on a torus, with torus-free constants}]\label{lem:a06A:torus}
Let $T$ be a torus $\eta\Z^d/\Gamma$ whose period lattice $\Gamma$ is contained in $M\Z^d$ (so the large blocks of the unit lattice tile $T$) and whose sup-systole is at least $3M$; let $\Omega\subset T$ be a union of large blocks ($\Omega=T$ included), $A$ regular on $\Omega$, and $e_k$ small. Then the theorem of \cite{Ba4} (\cite[Prop.~2.1, (2.24)--(2.26)]{Ba1}) holds on $T$ in its restricted form, with distances the quotient sup-distances and with the constants $\delta_0$, $c_0$, $R_0$ of \cite{Ba4}, which do not depend on $\Gamma$; for $\Omega=T$ it holds for all $x,x'\in T$. The same is true for tori of sup-systole less than $3M$ (but at least $M$) whose period lattice lies in $M\Z^d$. Consequently every estimate of \cref{thm:a06A:groupA} holds on $T$ with the constants of \cref{thm:a06A:groupA}.
\end{lemma}

\begin{proof}
The proof of \cite[Sect.~2]{Ba4} (pp.~574--580) is carried onto $T$ step by step; it uses the geometry of the lattice in the following places, and nowhere else.

(1) \emph{The local cubes.} Let $J$ index the vertices $Mj$ of the large-block partition of $T$ and $C_j\defeq Mj+[-M,M]^d$ (in unit-lattice units). Since $2M$ is less than the sup-systole, the projection $\eta\Z^d\to T$ is injective on each $C_j$; it preserves bonds, unit blocks, $L$-blocks and the contours of $Q_k$ inside blocks. So $\Pi_j\defeq\Omega\cap C_j$ \cite[p.~575]{Ba4} is a union of at most $2^d$ large blocks inside an embedded cube of side $2M$, a cube when $Mj$ is not a boundary point of $\Omega$, and the decomposition $A=A_0+A'$ on $C_j$ with the constant $A_0\defeq A(Mj)$ and the local field $A_j\defeq A_0+\theta_jA'$ of \cite[p.~575]{Ba4} is built from $A|_{C_j}$ alone, with $|A'|,|\partial^\eta A'|\le c'e_k^{\beta-1}$, $c'=dMc$. Hence \cite[Lemma~2.1]{Ba4} applies to every $G_k(\Pi_j,A_j)$ and \cite[Lemma~2.2]{Ba4} to the interior cubes (for $\Pi_j$ meeting $\partial\Omega$, $A_j\defeq A$ \cite[p.~575]{Ba4}), with their constants, which depend on $d$, $M$, $\alpha$, $c$ only; \cite[Lemma~2.4]{Ba4}, used in the proof of \cite[Lemma~2.2]{Ba4}, is \cref{thm:a06:A5} on parallelepipeds of the fine lattice.

(2) \emph{The partition of unity.} With $h\in C_0^\infty((-\frac23,\frac23))$, $h=1$ on $[-\frac13,\frac13]$, $\sum_{n\in\Z}h(\cdot-n)^2=1$, and $h_j(x)\defeq\prod_\mu h(x_\mu/M-j_\mu)$ \cite[p.~575]{Ba4}: because $\Gamma\subset M\Z^d$ the family $\{h_j\}_{j\in J}$ is well defined on $T$, $\sum_{j\in J}h_j^2=1$ on $T$ and $\supp h_j\subset C_j$. So the approximate inverse $G_0=\sum_jh_jG_k(\Pi_j,A_j)h_j$ \cite[(2.2)]{Ba4} and the identities \cite[(2.3)--(2.9)]{Ba4} hold on $T$ as on $\eta\Z^d$; for $\Omega=T$ the boundary of $\Omega$ is empty and the boundary term of \cite[(2.6)]{Ba4} vanishes.

(3) \emph{The remainder.} Each $j$ has at most $3^d-1$ neighbours and each point lies in at most $2^d$ supports, so the operator $R$ of \cite[(2.11)]{Ba4} has $\norm R_{2,2}\le c_dO(M^{-1})<1$ by \cite[Lemma~2.1]{Ba4}, as on \cite[p.~577]{Ba4}; the Neumann series \cite[(2.12)]{Ba4} converges and gives the walk representation \cite[(2.13)]{Ba4}, the walks being sequences of indices $j\in J$ with consecutive terms adjacent, and the count $2^d(3^d)^{n-1}2^d$ of walks of length $n$ \cite[p.~579]{Ba4} is unchanged.

(4) \emph{The metric.} Consecutive cubes of a walk have centres at quotient distance at most $M$, so a walk of length $n$ from a cube containing $x,x'$ to a cube containing points of $\supp f$ has $\dist(\{x,x'\},\supp f)\le(n+2)M$, i.e.\ $n\ge M^{-1}\dist(\{x,x'\},\supp f)-2$, which is the restriction used in \cite[(2.22)]{Ba4}; with (3), the geometric sum \cite[(2.22)]{Ba4} gives $\delta_0$ and $c_0$ as on $\eta\Z^d$, once $M$ is fixed with $3^dc_{2.1}O(1)M^{-1}\le\mathrm e^{-1}$, $c_{2.1}$ the constant of \cite[Lemma~2.1, (2.15)]{Ba4} (not the $c_2$ of \cref{lem:a06A:blockinv}). The definition \cite[(2.19)]{Ba4} of $R_0$ (the diameter of the union of the first $n_0$ cubes of a walk, plus $2M$) and the implication ``$\dist(\{x,x'\},\Omega^c)\ge R_0$, hence the first $n_0$ cubes $\Pi_{\omega_i}$ of every walk from $x,x'$ are cubes contained in $\Omega$'' use only the quotient metric and (1); on these cubes $A_j$ is constant near $\partial\Pi_j$ and \cite[Lemma~2.2]{Ba4} applies, on the others only \cite[Lemma~2.1]{Ba4} is used, as in \cite[(2.20)--(2.21)]{Ba4}. The cancellation of the terms with interior cubes in the two representations of $\delta G_k(\Omega,\Omega_0,A)$ \cite[p.~579]{Ba4} uses (1)--(2) only. If $\Omega=T$, every $\Pi_j=C_j$ is a cube with $A_j$ constant near its boundary, so no restriction on $x,x'$ is needed \cite[p.~611]{Ba1}.

(5) \emph{Small tori.} If the sup-systole of $T=\eta\Z^d/\Gamma$ is less than $3M$, let $\tilde T\defeq\eta\Z^d/3\Gamma$, a covering torus of systole at least $3M$ with deck group $(\Z/3)^d$. The operators $-\Delta^{\eta,N}_{A,\Omega}+m_k^2+a_kP_k(A)$ lifted to $\tilde T$ commute with the deck translations and map deck-invariant functions to deck-invariant functions, so for $f$ on $\Omega$ the lift of $G_k(\Omega,A)f$ is $G_k(\tilde\Omega,\tilde A)\tilde f$, with $\tilde\Omega,\tilde A,\tilde f$ the lifts; $\norm{\tilde f}_\infty=\norm f_\infty$, and the quotient distances in $\tilde T$ from a lift of $x$ to $\supp\tilde f$ and to $\tilde\Omega^c$ equal those in $T$. So the bounds on $\tilde T$, proved in (1)--(4), descend to $T$ with the same constants. The same applies to $\delta G_k$.

No other constant of \cite[Sect.~2]{Ba4} involves the global geometry. Finally, \cref{lem:a06A:cube,lem:a06A:coer,lem:a06A:CT,lem:a06A:blockinv,lem:a06A:resdiff,lem:a06A:CLdiff}, the first bound of \cref{prop:a06A:Ba1P22} and \cref{lem:a06A:polylog} are proved above on an arbitrary union of blocks of an arbitrary torus (their only global ingredients are the sums $S(\theta)$ and $S_L(\theta)$ over the block lattice, bounded by their values on $\Z^d$), and \cref{lem:a06A:twosided,lem:a06A:domdiff} and the second bound of \cref{prop:a06A:Ba1P22} use the printed theorem only in the form just established. Hence \cref{thm:a06A:groupA} holds on $T$ with the same constants.
\end{proof}



